\documentclass[11pt,a4paper]{article}

\usepackage[utf8]{inputenc}
\usepackage[T1]{fontenc}
\usepackage[polish]{babel}
\usepackage{lmodern}
\usepackage{geometry}
\geometry{margin=27mm}
\usepackage[hidelinks]{hyperref}

\title{Wpływ obustronnej stymulacji słuchowej na regulację\\
\large Projekt pilotażu dotyczącego czasu wznowienia zadania po zakłóceniu}
\author{Mateusz Pomianek}
\date{Prospektywny szkic badania empirycznego, wersja 0.4.0}

\begin{document}
\maketitle

\begin{abstract}
Celem opracowania jest przedstawienie projektu pilotażowego badania wpływu naprzemiennej stymulacji słuchowej na czas wznowienia zadania po zakłóceniu. Proponowany schemat wewnątrzosobniczy przewiduje porównanie naprzemiennej prezentacji dźwięku w obu kanałach słuchowych z warunkiem kontrolnym o porównywalnej głośności oraz losowaną lub równoważoną kolejnością ekspozycji. Kandydatem na główny punkt końcowy jest czas od ustalonego zdarzenia kończącego zakłócenie do wznowienia czynności. Odrębnie planuje się ocenę wykonalności procedury, tolerancji bodźców i kompletności pomiarów. Przed rozpoczęciem badania należy określić populację, parametry akustyczne, kryteria wznowienia, sposób traktowania prób bez zdarzenia i plan analizy oraz rozstrzygnąć wymagania etyczne. Szkic nie zawiera danych ani wyników empirycznych; czas wznowienia nie jest bezpośrednią miarą regulacji emocjonalnej ani skuteczności terapeutycznej.
\end{abstract}

\noindent\textbf{Słowa kluczowe:} naprzemienna stymulacja słuchowa; wznowienie zadania; zakłócenie zadania; regulacja zachowania; badanie pilotażowe; pomiar behawioralny.

\medskip
\noindent\textit{Rodzaj opracowania:} prospektywny konspekt badania empirycznego (pilotaż); nie jest raportem z przeprowadzonego eksperymentu.

\section{Wprowadzenie}
Przerwanie wykonywanej czynności wymaga ponownego ukierunkowania uwagi, odtworzenia celu i wznowienia sekwencji działań. Czas potrzebny do podjęcia zadania po zakłóceniu może opisywać widoczny aspekt powrotu do aktywności. Nie identyfikuje jednak samodzielnie mechanizmu samoregulacji. Na wynik mogą wpływać pamięć robocza, trudność i atrakcyjność zadania, natężenie reakcji orientacyjnej, zrozumienie polecenia, motywacja oraz warunki środowiskowe. Rozróżnienie tych składników jest warunkiem ostrożnej interpretacji zależności między bodźcem słuchowym a zachowaniem.

W planowanym pilotażu porównane mają zostać dwa sposoby prezentacji dźwięku: naprzemienne podawanie bodźca do lewego i prawego kanału słuchowego oraz warunek kontrolny o porównywalnej głośności. Termin „obustronna stymulacja słuchowa” określa cechę techniczną ekspozycji, nie dowodzi odrębnego mechanizmu neuronalnego ani efektu terapeutycznego. Nawet przy podobnym poziomie dźwięku oba warunki mogą różnić się percepcyjną wyrazistością, zmiennością przestrzenną, przewidywalnością i stopniem angażowania uwagi. Różnice te należy uwzględnić przy projektowaniu kontroli.

Projekt koncentruje się na pytaniu, czy sposób prezentacji dźwięku może wiązać się z czasem wznowienia zadania po jednolicie zdefiniowanym zakłóceniu. Jest to pytanie o wynik behawioralny, a nie o redukcję objawów ADHD, poprawę regulacji emocjonalnej czy skuteczność interwencji klinicznej. Populacja badana pozostaje do określenia; identyfikator projektu nie przesądza o rozpoznaniu u uczestników. W pierwszej kolejności konieczna jest ocena wykonalności pomiaru i tolerancji warunków akustycznych.

\section{Podstawy teoretyczne i zakres przeglądu literatury}
\subsection{Powrót do zadania a konstrukty psychologiczne}
Wznowienie zadania po zakłóceniu można zoperacjonalizować jako pierwsze obserwowalne działanie spełniające uprzednio ustalone kryterium kontynuowania aktywności. Taki wskaźnik nie jest tożsamy z odzyskaniem pełnej wydajności, stabilizacją uwagi ani regulacją emocji. Uczestnik może szybko wykonać pojedynczy ruch, lecz nadal nie odtworzyć wcześniejszej strategii działania. Odwrotnie, dłuższa przerwa może wynikać z przemyślanego wyboru lub ostrożności, a nie z trudności regulacyjnych.

Aby uzasadnić interpretację czasu wznowienia, potrzebne jest określenie cech zadania, momentu zakłócenia, kryterium ponownego podjęcia działania oraz granicznego czasu obserwacji. Należy też ocenić, czy wynik zależy od kolejnych prób i znajomości zadania. Ewentualne rozszerzenie interpretacji na regulację emocji wymagałoby odrębnej walidowanej miary, której nie przewidziano jeszcze w zamkniętym protokole.

\subsection{Możliwe mechanizmy oddziaływania stymulacji słuchowej}
Pierwszym rozważanym mechanizmem jest kierowanie uwagi przez zmieniającą się lokalizację bodźca. Naprzemienność kanałów może zwiększać liczbę zmian orientacji uwagowej, lecz może także stanowić dodatkowe rozproszenie. Drugi mechanizm dotyczy ogólnego pobudzenia. Bodziec słuchowy może modyfikować gotowość do działania niezależnie od lateralizacji, dlatego zrównanie głośności jest konieczne, ale niewystarczające do oddzielenia wszystkich składowych ekspozycji.

Trzecie wyjaśnienie dotyczy przewidywalności i rytmiczności prezentacji. Regularny sygnał może organizować moment działania, podczas gdy sygnał zmienny może odwracać uwagę. Czwarty mechanizm obejmuje zróżnicowaną tolerancję sensoryczną: identyczne parametry fizyczne mogą być odmiennie odbierane przez uczestników. Piątym źródłem różnic mogą być oczekiwania wobec warunku eksperymentalnego oraz efekty kolejności, ćwiczenia i zmęczenia. Są to konkurencyjne wyjaśnienia wymagające kontroli, a nie potwierdzone mechanizmy działania badanej procedury.

\subsection{Problem identyfikowalności efektu}
Porównanie tylko dwóch warunków umożliwia oszacowanie ich różnicy, ale nie musi pozwalać na przypisanie tej różnicy wyłącznie naprzemienności bodźców. Jeżeli warunki różnią się jednocześnie rytmem, przestrzenną lokalizacją i przewidywalnością, interpretacja przyczynowa będzie dotyczyć całego pakietu cech ekspozycji. W przyszłym protokole należy albo wyrównać te cechy, albo wyraźnie ograniczyć zakres wnioskowania. Nie ma obecnie podstaw do deklarowania separacji poszczególnych mechanizmów.

\subsection{Luka badawcza i zadanie weryfikacji literatury}
Roboczą luką metodologiczną jest potrzeba sprawdzenia, czy krótki, standaryzowany pomiar czasu wznowienia umożliwi wiarygodne porównanie warunków słuchowych, uwzględniające próby bez wznowienia i indywidualną tolerancję dźwięku. Nie jest to stwierdzenie, że analogiczne badania nie istnieją. Ocena nowości wymaga udokumentowanego przeglądu piśmiennictwa oraz porównania populacji, konfiguracji bodźców, warunków kontrolnych, kryteriów wyniku i ryzyka błędu systematycznego.

Przed nadaniem tekstowi postaci artykułu z przeglądem literatury należy zweryfikować publikacje dotyczące przerywania i wznawiania zadań, interwencji z obustronnymi bodźcami słuchowymi, orientowania uwagi, wrażliwości sensorycznej oraz metod analizy czasów do zdarzenia. Dla każdego twierdzenia empirycznego konieczne będzie wskazanie rzeczywiście sprawdzonego źródła. W niniejszym etapie nie wprowadzono niezweryfikowanych cytowań ani twierdzeń o potwierdzonej skuteczności.

\section{Cel, pytania badawcze i status hipotez}
\subsection{Cel główny i cele pomocnicze}
Głównym celem planowanego pilotażu jest ocena wykonalności porównania naprzemiennej stymulacji słuchowej z kontrolą akustyczną oraz zbadanie możliwości oszacowania różnicy czasu wznowienia zadania po standardowym zakłóceniu. Cele pomocnicze obejmują ocenę kompletności rejestracji, częstości prób bez wznowienia w oknie obserwacji, akceptowalności ekspozycji i zmienności pomiaru między próbami oraz uczestnikami. Ustalenie liczbowych progów wykonalności wymaga osobnej decyzji przed rekrutacją.

\subsection{Pytania badawcze}
Podstawowe pytanie dotyczy tego, czy przy jednakowo zdefiniowanym zadaniu i zakłóceniu rozkład czasu wznowienia różni się między dwoma warunkami akustycznymi. Drugie pytanie odnosi się do tego, czy liczba i jakość zarejestrowanych prób umożliwią wiarygodne oszacowanie różnicy w pilotażu. Trzecie pytanie dotyczy tego, jak często pomiar zostanie zakończony bez wznowienia zadania i czy będzie to wymagało modelowania czasu do zdarzenia. Pytania dodatkowe obejmują zależność wyników od kolejności warunków i zgłaszanej tolerancji dźwięku.

\subsection{Hipotezy możliwe do ustalenia przed badaniem}
Na poziomie projektu możliwe jest sformułowanie niedyrektywnego kontrastu, w którym porównuje się czas wznowienia w warunku naprzemiennej stymulacji słuchowej z czasem w warunku kontrolnym. Nie przesądza się znaku różnicy. Samo zdefiniowanie kontrastu nie oznacza jeszcze prerejestracji hipotezy konfirmacyjnej ani przesądzenia, że pilotaż ma wystarczającą moc do jej testowania.

Hipoteza kierunkowa, przewidywana wielkość efektu, estymand, pierwszorzędowy punkt końcowy, zasady obsługi cenzorowania i próg decyzyjny wymagają uzasadnienia literaturowego oraz zapisania przed dostępem do danych. Do tego czasu nie należy oznaczać jakiejkolwiek hipotezy jako zatwierdzonej konfirmacyjnie. Cele pilotażowe dotyczą przede wszystkim oszacowania wykonalności i niepewności, a nie rozstrzygnięcia o klinicznej skuteczności procedury.

\subsection{Pytania eksploracyjne}
Eksploracyjnie można rozważyć związek tolerancji bodźców z czasem wznowienia, różnice pomiędzy kolejnymi próbami, wpływ kolejności warunków oraz ewentualną heterogeniczność efektu między uczestnikami. Analizy te powinny być wyraźnie odróżnione od wcześniej ustalonych porównań głównych. Nie wolno przedstawiać odkrytych po analizie zależności jako hipotez sformułowanych przed badaniem.

\subsection{Warunki interpretacji}
Wnioskowanie przyczynowe o sposobie prezentacji dźwięku będzie uzależnione od poprawnego przydziału warunków, kontroli cech akustycznych, zachowania procedury i braku nierozpoznanych efektów przenoszenia. Ewentualny krótszy czas wznowienia może mieć wiele wyjaśnień i sam w sobie nie dowodzi poprawy regulacji emocjonalnej, funkcji wykonawczych ani objawów ADHD. Interpretacja mechanistyczna i kliniczna wymagałaby niezależnego programu walidacyjnego.

\section{Projekt badania i metody}
\subsection{Schemat i jednostka obserwacji}
Rozważany jest pilotaż wewnątrzosobniczy z losowaną lub równoważoną kolejnością warunków. Każdy uczestnik wykonywałby zadanie obejmujące standardowe zakłócenie i późniejszą możliwość wznowienia aktywności. Należy rozstrzygnąć liczbę sesji i prób, odstępy między nimi, sposób ograniczania efektu uczenia i przenoszenia oraz z góry określić jednostkę analizy.

\subsection{Populacja i dobór próby}
Populacja docelowa nie została jeszcze zdefiniowana. Przed rekrutacją konieczne jest określenie przedziału wieku, kryteriów włączenia i wyłączenia, sposobu ustalania charakterystyki neurorozwojowej, planu doboru liczebności i dostępności uczestników. Sam identyfikator ADHD w nazwie projektu nie stanowi dowodu rozpoznania u przyszłych osób badanych.

\subsection{Bodźce i warunek kontrolny}
Warunek eksperymentalny zakłada naprzemienną prezentację dźwięku w kanałach słuchowych. Warunek kontrolny ma zachowywać porównywalny poziom głośności i czas prezentacji. Protokół powinien określić sposób kalibracji, rodzaj dźwięku, częstotliwość zmian kanału, czas ekspozycji, sposób odtwarzania oraz akceptowalne warunki otoczenia. Bez tych decyzji porównanie nie ma jeszcze wystarczającej jednoznaczności.

\subsection{Przebieg pojedynczej próby}
Planowana sekwencja obejmuje instrukcję zadania, wykonanie czynności początkowej, standardowe zakłócenie, ekspozycję na przypisany warunek akustyczny oraz obserwację wznowienia zadania. Konieczne jest precyzyjne określenie chwili rozpoczęcia pomiaru i operacyjnego kryterium wznowienia. Obserwator lub system rejestracji powinien stosować jednakowe reguły w każdym warunku.

\subsection{Zmienne i jakość pomiaru}
Kandydatem na główny punkt końcowy jest czas od zdefiniowanego zdarzenia kończącego zakłócenie do pierwszego spełnienia kryterium wznowienia. Należy z góry określić dopuszczalne okno obserwacji i sposób oznaczania prób bez wznowienia. Dodatkowymi zmiennymi mogą być kolejność warunków, tolerancja dźwięku i poziom hałasu, o ile ich wiarygodny pomiar zostanie przewidziany w protokole. Rejestr powinien umożliwiać identyfikację braków danych i zdarzeń zakłócających, bez ujawniania danych osobowych.

\subsection{Bezpieczeństwo, etyka i ochrona danych}
Przed rozpoczęciem udziału ludzi należy ustalić wymaganą ścieżkę oceny etycznej, świadomą zgodę lub zgodę opiekuna, zasady rezygnacji bez konsekwencji oraz procedury postępowania przy dyskomforcie sensorycznym. Szczególnej kontroli wymagają poziom głośności i nadwrażliwość słuchowa. Nie zakłada się uzyskania zgody etycznej ani zatwierdzenia urządzeń.

\section{Plan analizy}
\subsection{Estymand i porównanie warunków}
Proponowanym przedmiotem estymacji jest różnica czasu wznowienia między warunkami w obrębie uczestnika, w populacji, która zostanie uprzednio określona. Definicja podsumowania różnic, traktowanie prób niewznowionych i docelowa precyzja oszacowania wymagają decyzji przed analizą.

\subsection{Model statystyczny}
Można rozważyć model mieszany z efektem warunku, kolejności i uzasadnionych współzmiennych oraz efektem losowym uczestnika. Dobór rodziny rozkładu i funkcji łączącej powinien wynikać z charakteru czasów reakcji i sposobu cenzorowania. Jeśli wystąpią liczne próby bez wznowienia, analizę czasu do zdarzenia należy rozważyć zamiast arbitralnego przypisywania im maksymalnego czasu.

\subsection{Dane brakujące i analizy wrażliwości}
Należy odróżnić brak rejestracji od niewystąpienia zdarzenia w oknie obserwacji. Przed dostępem do wyników trzeba określić kryteria wyłączenia prób, plan kontroli odstających obserwacji, wpływu kolejności i tolerancji sensorycznej oraz raportowania niepewności oszacowań. Pilotaż nie powinien być interpretowany jako potwierdzający skuteczność kliniczną wyłącznie na podstawie istotności statystycznej.

\section{Przewidywany sposób prezentacji wyników}
Po uzyskaniu zweryfikowanych danych raport empiryczny powinien przedstawiać liczebność i przepływ uczestników, wykonalność procedury, kompletność pomiarów, rozkłady czasu wznowienia, różnice między warunkami i przedziały niepewności. Obecny tekst stanowi wyłącznie plan raportowania, a nie sekcję z wynikami przeprowadzonego badania.

\section{Dyskusja i ograniczenia planu}
Przyszła interpretacja powinna rozważyć, czy zaobserwowana zmiana odzwierciedla specyficzny wpływ naprzemienności, niespecyficzne pobudzenie, orientowanie uwagi, efekt oczekiwań lub różnice w tolerancji bodźców. Istotne są także efekty kolejności, uczenie się zadania i ograniczona reprezentatywność próby pilotażowej.

\subsection{Granice wnioskowania}
Czas wznowienia nie jest równoważny bezpośredniemu pomiarowi regulacji emocjonalnej. Bez dodatkowych pomiarów i kontroli nie będzie możliwe wnioskowanie o mechanizmie neuronalnym, trwałości efektu ani skuteczności terapeutycznej. Projekt musi także odróżniać ocenę wykonalności procedury od testowania efektu naukowego.

\section{Wnioski}
Opracowano strukturę prospektywnego badania pilotażowego, wskazując główny wskaźnik behawioralny, projekt porównania i wymagania analityczne. Kolejne kroki obejmują krytyczny przegląd literatury, domknięcie protokołu, rozstrzygnięcia etyczne i dopiero następnie ewentualne gromadzenie danych. Na tym etapie nie ma podstaw do formułowania wniosków empirycznych.

\section*{Informacje o autorstwie i oświadczenia}
\noindent\textbf{Afiliacja autora.} Do uzupełnienia po potwierdzeniu aktualnej afiliacji właściwej dla tej pracy. Nie przypisano jednostki organizacyjnej bez weryfikacji.

\medskip
\noindent\textbf{ORCID.} Identyfikator autora wymaga potwierdzenia i uzupełnienia. Nie podano niezweryfikowanego numeru.

\medskip
\noindent\textbf{Wkład autorów (CRediT).} Mateusz Pomianek jest wskazany jako autor szkicu. Szczegółowe role w taksonomii CRediT, w tym conceptualization, methodology, formal analysis, investigation, data curation, writing -- original draft oraz writing -- review \& editing, należy przypisać wyłącznie po potwierdzeniu rzeczywistego wkładu. Nie wskazano niepotwierdzonych współautorów ani wykonanych czynności badawczych.

\medskip
\noindent\textbf{Finansowanie.} Informacja o źródłach finansowania wymaga potwierdzenia przed zgłoszeniem tekstu. Nie deklaruje się ani przyznania środków, ani ich braku.

\medskip
\noindent\textbf{Zgoda etyczna i zgoda uczestników.} Badanie opisano wyłącznie jako plan. Nie przedstawiono decyzji komisji etycznej, rekrutacji ani zgód uczestników. Wymagane oświadczenia zostaną sformułowane zgodnie z rzeczywistym przebiegiem i wymaganiami właściwej instytucji.

\medskip
\noindent\textbf{Dostępność danych i materiałów.} W ramach niniejszego szkicu nie pozyskano materiału empirycznego. Zasady udostępniania przyszłych danych i protokołu wymagają odrębnej decyzji, uwzględniającej ochronę prywatności i granice repozytorium kanonicznego.

\medskip
\noindent\textbf{Konflikt interesów.} Oświadczenie autora wymaga weryfikacji przed zgłoszeniem. Braku konfliktu interesów nie przyjęto domyślnie.

\medskip
\noindent\textbf{Rejestracja i identyfikatory publikacji.} Nie potwierdzono prerejestracji, rejestracji protokołu ani identyfikatora DOI. Pola należy uzupełnić wyłącznie na podstawie dokumentów źródłowych.

\end{document}